\chapter*{Ringraziamenti}
\addcontentsline{toc}{chapter}{Ringraziamenti}
Alla fine di un percorso tanto importante, che dalle scuole elementari ad oggi ha plasmato la mia persona, è naturale ringraziare tutti coloro che, in modi diversi, hanno fatto parte di questa formazione.

In primo luogo i miei genitori: grazie alla mia mamma che è riuscita a trasmettermi la capacità di entusiasmarmi nello scoprire cose nuove e grazie al mio papà che ha sempre fomentato la mia curiosità e incoraggiato la mia voglia di sapere. Solo grazie a loro so cos'è la disciplina e quanto sia fondamentale per raggiungere ogni obiettivo. Grazie al loro sostegno, morale e materiale, e alla fiducia che hanno avuto nelle mie capacità ho potuto seguire liberamente i miei sogni ed arrivare fino a questo punto. 

Grazie alle mie sorelle, Lucrezia da una parte con la sua genialità e il trasporto in quello che fa e Ludovica dall'altra con la sua razionalità e sicurezza che la rendono un punto di riferimento per me e per tutti nonostante la differenza d'età. Loro sanno consolarmi quando sono triste, sgridarmi quando sbaglio, spronarmi quando non ho voglia. Sono grata di avervi avute come compagne di giochi quando eravamo bambine e come confidenti, amiche e punti di riferimento adesso che siamo cresciute e per sempre. 

Grazie ai miei nonni: a chi purtroppo non c'è più e a chi veglia ancora quotidianamente su di me. In modi diversi mi hanno tutti accompagnato e aiutato in questi anni. 

Grazie a mia zia Elena che è stata la mia quotidianità in questi sei anni a Padova e che mi ha sostenuto molto più di quanto lei creda in tutte le tappe di questo percorso. Grazie a mia zia Anna e alle mie cugine, Giorgia e Martina, sempre pronte ad aiutarmi in qualsiasi imprevisto della mia vita universitaria.

Grazie a Francesco, che ogni giorno mi supporta incondizionatamente, che mi vuole bene quando sono felice e ancora di più quando sono triste o arrabbiata, che mi sprona, che mi incoraggia, che porta (tanta) pazienza, che mi fa ridere anche quando non ne ho voglia, che mi calma quando non ho più controllo, che mi appoggia in ogni scelta e che spero farà tutte queste cose ancora per tanto tempo.

Ringrazio Irene e Laura per essere le migliori amiche che potessi desiderare, per sopportarmi ancora dopo più di dieci anni, per i consigli che sanno darmi dopo avermi ascoltato, per la pazienza che portano quando mi incastro nei miei pensieri, per tutte le serate che sono state e per tutte quelle che verranno.

Grazie ad Elisa e Federica perché quando ci vediamo sembra che il tempo non sia mai trascorso. Grazie ad Alessandro, Luca, Nicola e Giovanni perché con voi avrò sempre diciassette anni. Grazie alla III C che dopo anni continua ad essere un punto di riferimento nonostante le vite molto diverse che conduciamo e qualche centinaia di chilometri che a volte ci separa.

Grazie ad Anna, sempre ad un messaggio di distanza nonostante i chilometri e il tempo che passa.

Grazie a Martina, che ha reso ogni giornata di studio e di laboratorio una fonte infinita di risate, divertimento e ricordi indelebili.

Grazie a Veronica, Caterina, Emma, Giulia, Cristian e Luca: chi solo durante la Triennale e chi anche negli ultimi anni della Magistrale, hanno tutti accompagnato le mie giornate universitarie tra discorsi di incoraggiamento, pranzi condivisi e sessioni interminabili di studio.

Grazie ad Elena per aver condiviso con me gran parte del percorso della Magistrale, grazie a Sergio per i pomeriggi passati a sbagliare tutto in laboratorio e grazie a Michelangelo, Tommaso e Marco per tutte le risate, per le pause caffè e per avermi fatto scommettere sui miei esami: ne è valsa la pena.